\scanpage{228}
% continuation, reduction about a point
point $s_0\in S$ on est ramené \unsure{comme|ainsi|encore} au cas où $a(s_0)$ est proche de $b(s_0)$, de sorte qu'on peut \deleted{\unsure{se pla}} se placer dans un ouvert de $X$ assez petit pour qu'il y ait une \unsure{connection|connexion}. Et on gagne.

% alternate formulation
On peut aussi noter, \deleted{\unsure{qu'on|quitte}} pour prouver la \deleted{formule} nullité de (3) comme conséquence de (7), que l'on a
\[
 \sum D_{i\,s}\int_{Z_s}\omega_i
 =\int_{Z_s}\bigl(\sum\overline D_i\,\omega'_i\bigr)_s
 \tag*{\deleted{\unsure{(15)}}}
\]
où le deuxième mb a un sens, car $\int_{Z_s}\overline\omega_s=0$ si $\overline\omega_s$ est un bord, \unsure{or|et} $\bigl(\sum\overline D_i\,\omega'_i\bigr)_s$ est une forme définie modulo un bord. La formule envisagée \deleted{\unsure{équivaut à}} \inserted{revient à}
\[
 \boxed{D_s\int_{Z_s}\omega'_s=\int_{Z_s}(\overset{\deleted{\unsure{-}}}{D}\omega')_s}
 \tag*{(15)}
\]
qui donne la propriété transcendante essentielle de la connexion intégrable canonique sur $\Phi$. Nous n'allons pas prouver (15) ici, et l'admettrons, mais notons que
% archive sheet number: 227
